% Nini — guide complet d'utilisation
% Compiler avec XeLaTeX (polices Inter et DejaVu installées), deux passes pour le sommaire :
%   xelatex nini-guide-complet.tex && xelatex nini-guide-complet.tex
% Les captures d'écran (dossier images/) utilisent uniquement les données d'exemple de l'app.
\documentclass[11pt,a4paper]{article}

\usepackage[french]{babel}
\usepackage{fontspec}
\setmainfont{Inter}[
  UprightFont = *-Regular, BoldFont = *-SemiBold,
  ItalicFont = *-Italic, BoldItalicFont = *-SemiBoldItalic, Extension = .otf ]
\newfontfamily\displayfont{InterDisplay}[UprightFont = *-Bold, BoldFont = *-Bold, Extension = .otf]
\setmonofont{DejaVu Sans Mono}[Scale=0.85]
\newfontfamily\symb{DejaVu Sans}

\usepackage[a4paper, margin=21mm, top=19mm, bottom=23mm]{geometry}
\usepackage{xcolor}
\usepackage{graphicx}
\usepackage{microtype}
\usepackage{enumitem}
\usepackage{tabularx}
\usepackage{colortbl}
\usepackage{titlesec}
\usepackage{titletoc}
\usepackage{needspace}
\usepackage{fancyhdr}
\usepackage{tikz}
\usetikzlibrary{babel}
\usepackage[most]{tcolorbox}
\usepackage[hidelinks, pdftitle={Nini, guide complet}, pdfauthor={Nini}]{hyperref}

% ---------- Couleurs (celles de l'app) ----------
\definecolor{accent}{HTML}{3355FF}
\definecolor{accentsoft}{HTML}{EEF1FF}
\definecolor{ink}{HTML}{17171C}
\definecolor{ink2}{HTML}{52525C}
\definecolor{line}{HTML}{E2E3E9}
\definecolor{soft}{HTML}{F5F6F9}
\definecolor{ok}{HTML}{1D9A61}
\definecolor{oksoft}{HTML}{E9F7F0}
\definecolor{no}{HTML}{D93F45}
\definecolor{nosoft}{HTML}{FDEEEE}
\definecolor{amber}{HTML}{B8780A}
\definecolor{ambersoft}{HTML}{FFF5E3}
\definecolor{pink}{HTML}{E0569A}
\color{ink}

% ---------- Mise en page ----------
\setlength{\parindent}{0pt}
\setlength{\parskip}{6pt}
\linespread{1.15}
\setcounter{secnumdepth}{2}
\setcounter{tocdepth}{1}
\titleformat{\section}{\needspace{12\baselineskip}\displayfont\fontsize{19}{23}\selectfont}
  {\color{accent}\thesection}{10pt}{}
\titlespacing*{\section}{0pt}{26pt}{10pt}
\titleformat{\subsection}{\needspace{7\baselineskip}\bfseries\large}{}{0pt}{}
\titlespacing*{\subsection}{0pt}{16pt}{6pt}
\titlecontents{section}[1.8em]{\addvspace{1pt}}{\contentslabel{1.8em}}{}{\titlerule*[0.6pc]{.}\contentspage}
\titlecontents{subsection}[3.6em]{\small\color{ink2}}{\contentslabel{1.8em}}{}{\titlerule*[0.6pc]{.}\contentspage}
\pagestyle{fancy}
\fancyhf{}
\renewcommand{\headrulewidth}{0pt}
\fancyfoot[L]{\small\color{ink2}Nini · guide complet}
\fancyfoot[R]{\small\color{ink2}\thepage}

% ---------- Composants ----------
\newcommand{\muted}[1]{{\color{ink2}#1}}
\newcommand{\btn}[1]{\textbf{#1}}
\newcommand{\kbd}[1]{\tcbox[on line, boxsep=0pt, left=3pt, right=3pt, top=1.5pt, bottom=1.5pt,
  colback=soft, colframe=line, boxrule=0.5pt, arc=2pt]{\ttfamily\small #1}}
\newcommand{\pastille}[1]{\tcbox[on line, boxrule=0pt, colback=accentsoft, colframe=accentsoft, arc=8pt,
  boxsep=0pt, left=4.5pt, right=4.5pt, top=3pt, bottom=3pt]{\makebox[0.62em]{\small\bfseries\color{accent}#1}}}
\newlist{etapes}{enumerate}{1}
\setlist[etapes]{label=\pastille{\arabic*}, leftmargin=28pt, labelsep=8pt, itemsep=4pt, topsep=4pt}
\setlist[itemize]{leftmargin=16pt, itemsep=2pt, topsep=3pt, label={\color{accent}\textbullet}}

\newcommand{\capture}[3][0.9\linewidth]{% largeur, fichier, légende (gardées ensemble)
  \par\medskip\noindent\begin{minipage}{\linewidth}\centering
  \begin{tcolorbox}[enhanced, width=#1, boxrule=0.5pt, colframe=line, colback=white, arc=5pt,
    left=0pt, right=0pt, top=0pt, bottom=0pt, boxsep=0pt, before skip=0pt, after skip=3pt]
    \includegraphics[width=\linewidth]{images/#2}
  \end{tcolorbox}
  {\footnotesize\muted{#3}}
  \end{minipage}\par\medskip}

\newtcolorbox{astuce}{enhanced, colback=accentsoft, colframe=accentsoft, arc=7pt, boxrule=0pt,
  left=11pt, right=11pt, top=7pt, bottom=7pt, before skip=8pt, after skip=8pt,
  fontupper=\small, before upper={\textbf{\color{accent}Astuce\enspace}}}
\newtcolorbox{attention}{enhanced, colback=ambersoft, colframe=ambersoft, arc=7pt, boxrule=0pt,
  left=11pt, right=11pt, top=7pt, bottom=7pt, before skip=8pt, after skip=8pt,
  fontupper=\small, before upper={\textbf{\color{amber}Attention\enspace}}}
\newtcolorbox{securite}{enhanced, colback=oksoft, colframe=oksoft, arc=7pt, boxrule=0pt,
  left=11pt, right=11pt, top=8pt, bottom=8pt, before skip=8pt, after skip=8pt}

% Tableau à deux colonnes : terme / explication
\newenvironment{lexique}[1][0.28]
  {\par\renewcommand{\arraystretch}{1.45}\arrayrulecolor{line}%
   \tabularx{\linewidth}{@{}>{\bfseries\raggedright\arraybackslash}p{#1\linewidth}>{\raggedright\arraybackslash\leavevmode\color{ink2}}X@{}}\hline}
  {\endtabularx\par}

\begin{document}

% =====================================================================
% Couverture
% =====================================================================
\thispagestyle{empty}
\begin{tikzpicture}
  \node[fill=accent, rounded corners=7pt, minimum size=15mm, text=white, font=\displayfont\fontsize{25}{25}\selectfont] {N};
\end{tikzpicture}

\vspace{10pt}
{\displayfont\fontsize{36}{40}\selectfont Nini\par}
\vspace{2pt}
{\displayfont\fontsize{22}{26}\selectfont\color{ink2} Le guide complet\par}
\vspace{10pt}
{\large\muted{Tout ce qu’il faut savoir pour suivre tes dossiers, organiser ton mois et garder tes données en sécurité.}\par}

\vspace{14pt}
\begin{securite}
\textbf{À retenir avant tout.}\enspace Nini ne se connecte jamais à Internet : tes données restent sur ton appareil.
En contrepartie, \textbf{c’est toi qui fais les sauvegardes}, chaque vendredi.
\end{securite}

\vspace{6pt}
{\small\muted{Les captures d’écran de ce guide montrent des dossiers d’exemple, pas tes vrais clients.}}

\vspace{14pt}
{\displayfont\fontsize{15}{18}\selectfont Sommaire\par}
\vspace{2pt}
{\makeatletter\@starttoc{toc}\makeatother}

\newpage

% =====================================================================
\section{Nini en bref}
% =====================================================================
Nini est ton carnet de travail de collaboratrice comptable. Tu y retrouves :
\begin{itemize}
  \item \textbf{tes dossiers clients}, rangés par manager, avec le temps prévu pour chacun ;
  \item \textbf{les tâches du mois}, créées automatiquement pour chaque dossier, avec des étapes à cocher ;
  \item \textbf{ton planning}, qui répartit tes journées entre les dossiers ;
  \item \textbf{des rapports} sur ce qui est fait et ce qui reste à faire.
\end{itemize}

\subsection{Les mots à connaître}
\begin{lexique}
Dossier & Un client du cabinet (une société, une SCI…). Chaque dossier a un code court, par exemple \texttt{BOUM}. \\ \hline
Manager & La personne responsable du dossier. Les dossiers sont regroupés par manager. \\ \hline
Tâche & Un travail à faire : « Début de mois », « TVA \& révision »… \\ \hline
Étapes & La liste à cocher à l’intérieur d’une tâche : intégration banque, lettrage… \\ \hline
Période & Un moment du mois : début, milieu, TVA, fin. Chaque période a ses tâches. \\ \hline
Date de début & Le jour où la tâche apparaît dans « À traiter en ce moment ». \\ \hline
Échéance & La date à laquelle la tâche doit être finie. \\ \hline
Répartition & Le planning du mois : quel dossier, quel jour, combien de temps. \\ \hline
Créneau & Un bloc du planning : un dossier, un jour, une durée (1 h, 0,5 j…). \\ \hline
Sauvegarde & Un fichier qui contient toutes tes données, pour les protéger ou les transporter. \\ \hline
\end{lexique}

% =====================================================================
\section{Ouvrir Nini}
% =====================================================================
\subsection{Sur l’ordinateur}
\begin{etapes}
  \item Copie le fichier \btn{Nini.html} dans tes \btn{Documents}. Ne le déplace plus ensuite.
  \item Double-clique dessus : il s’ouvre dans \btn{Edge} ou \btn{Chrome}.
  \item Ajoute la page à tes favoris pour la retrouver en un clic.
\end{etapes}
\begin{attention}
Ouvre toujours Nini de la même façon : même fichier, même endroit, même navigateur.
Sinon, le navigateur ne retrouve pas tes données.
\end{attention}

\subsection{Sur l’iPad}
L’iPad ne sait pas ouvrir un fichier \texttt{.html}. On installe donc Nini comme une app, une seule fois (voir la partie \ref{sec:ipad}).

\subsection{La première fois}
\capture[0.68\linewidth]{00-bienvenue.jpg}{L’écran d’accueil quand Nini est vide.}
Trois possibilités :
\begin{itemize}
  \item \btn{Créer un premier dossier} pour partir de zéro ;
  \item \btn{Explorer avec des données d’exemple} pour découvrir l’app sans risque ;
  \item ou, le plus courant, \btn{importer tes dossiers} à partir du fichier de sauvegarde que tu as reçu :
  menu \btn{Réglages} → \btn{Restaurer…} → choisis le fichier \texttt{.json}.
\end{itemize}
\begin{astuce}
Après avoir essayé les données d’exemple, va dans \btn{Réglages → Tout effacer}, puis restaure ta sauvegarde.
\end{astuce}

% =====================================================================
\section{L’écran principal}
% =====================================================================
\capture{01-accueil.jpg}{À gauche, la barre latérale. À droite, la page ouverte (ici : Aujourd’hui).}

La \textbf{barre latérale}, à gauche, est toujours là :
\begin{lexique}[0.3]
Rechercher & Ouvre la palette de recherche (aussi avec \kbd{Ctrl}+\kbd{K}) : retrouve une tâche ou un dossier en tapant quelques lettres. \\ \hline
Nouvelle tâche & Crée une tâche (aussi avec la touche \kbd{N}). \\ \hline
Aujourd’hui & Ta journée : ce qui est prévu, ce qui est en retard, ce qui est à traiter. Le chiffre indique le nombre de tâches à traiter. \\ \hline
Toutes les tâches & La liste complète, avec des filtres. \\ \hline
Répartition & Le planning du mois. \\ \hline
Calendrier & Les échéances sur un calendrier. \\ \hline
Rapports & Le bilan de ce qui est fait et de ce qui reste. \\ \hline
Dossiers & La liste de tes clients. \\ \hline
Groupes par manager & En dessous, tes dossiers rangés par manager. Clique sur le nom d’un manager pour \textbf{déplier ou replier} sa liste. Le chiffre gris indique le nombre de dossiers, le chiffre rouge les retards. \\ \hline
Sauvegarde conseillée & Ce cadre jaune apparaît après 7 jours sans sauvegarde. Clique sur \btn{Sauvegarder}. \\ \hline
Chrono & Quand un chrono tourne, il s’affiche en bas, avec un bouton pour l’arrêter. \\ \hline
Réglages & Sauvegardes, routine du mois, sécurité. \\ \hline
\end{lexique}

% =====================================================================
\section{Aujourd’hui}
% =====================================================================
C’est la page à ouvrir chaque matin. De haut en bas :
\begin{lexique}[0.3]
Le résumé & Sous le titre : le nombre de tâches en retard, à traiter, pour aujourd’hui, dans les 7 jours, et chez le client. Clique dessus pour voir la liste. \\ \hline
Au programme aujourd’hui & Les dossiers prévus ce jour dans ta répartition, avec la durée et l’avancement des étapes (par exemple « 1 h · 2/8 »). Un clic ouvre directement la bonne tâche. La période en cours est rappelée juste au-dessus. \\ \hline
À reprendre & Les tâches que tu avais reportées et dont la date de reprise est arrivée. \btn{Reprendre} les remet en cours ; \btn{Reporter encore} rouvre la tâche pour choisir une autre date. \\ \hline
En retard & Les tâches dont l’échéance est passée. \\ \hline
Aujourd’hui & Les tâches à finir aujourd’hui, et un champ pour en ajouter une. \\ \hline
À traiter en ce moment & Toutes les tâches de la période en cours, groupées. Pour chaque groupe : la date de fin et le nombre de tâches terminées. \\ \hline
Les 7 prochains jours & Les échéances de la semaine, jour par jour. \\ \hline
En attente du client & Les tâches bloquées parce qu’il manque quelque chose au client, avec le nombre de jours depuis la relance. \\ \hline
\end{lexique}
\begin{astuce}
Le rond à gauche d’une tâche sert à la cocher. Si tu te trompes, clique sur \btn{Annuler} dans le message qui apparaît en bas.
\end{astuce}

% =====================================================================
\section{Travailler une tâche}
% =====================================================================
Clique sur une tâche : elle s’ouvre dans un panneau à droite. \textbf{Tout s’enregistre tout seul}, il n’y a pas de bouton « Enregistrer ».

\capture{02-panneau.jpg}{Le panneau d’une tâche : titre, informations, étapes.}

\subsection{Les informations}
\begin{lexique}
Statut & \textbf{À faire}, \textbf{En cours}, \textbf{En attente client} (il manque une pièce), \textbf{Terminé}. Passer en « En attente » démarre le compteur de jours depuis la relance. \\ \hline
Début & Le jour où la tâche apparaît dans « À traiter en ce moment ». \\ \hline
Échéance & La date de fin. À droite, Nini indique « dans 3 j », « demain » ou « en retard ». \\ \hline
Dossier & Le client concerné. Choisis \btn{Aucun — à lier plus tard} si tu ne sais pas encore : tu pourras le choisir ensuite. \\ \hline
Priorité & Urgente, Haute, Normale, Basse. Les tâches urgentes ont un rond rouge, les hautes un rond orange. \\ \hline
Catégorie & Saisie, TVA, Révision… Elle sert aux filtres et aux rapports. \\ \hline
Estimation & Le temps prévu, en heures. \\ \hline
Récurrence & Chaque semaine, mois, trimestre ou année : quand tu termines la tâche, Nini crée la suivante toute seule. \\ \hline
\end{lexique}
Sous les informations, une zone libre pour tes \textbf{notes} : pièces manquantes, consignes…

\subsection{Les étapes}
Coche chaque étape quand elle est faite. La barre bleue montre l’avancement.
Pour ajouter une étape, écris dans \btn{+ Ajouter une étape} puis appuie sur \kbd{Entrée}.
Pour en supprimer une, passe la souris dessus et clique sur la croix.

\capture{03-panneau-suite.jpg}{Plus bas dans le panneau : le temps passé et l’activité.}

\subsection{Le temps passé}
Deux façons de noter ton temps :
\begin{itemize}
  \item \textbf{le chrono} : bouton \btn{Chrono} en haut du panneau. Il continue même si tu changes de page. Arrête-le avec le même bouton, ou en bas de la barre latérale : le temps s’ajoute tout seul ;
  \item \textbf{à la main} : choisis la date, écris la durée (\texttt{1h30}, \texttt{45}, \texttt{0:30}…), une note si tu veux, puis \btn{Ajouter}.
\end{itemize}

\subsection{L’activité}
Écris un commentaire (appel client, relance…) puis \btn{Publier}. Nini note aussi tout seul les changements importants : statut, échéance, priorité, temps ajouté. C’est l’historique de la tâche.

\subsection{Les boutons du haut}
\begin{lexique}
Chrono & Démarre ou arrête le chrono de cette tâche. \\ \hline
Reporter & Met la tâche de côté jusqu’à une date de reprise (voir ci-dessous). \\ \hline
Dupliquer & Crée une copie de la tâche, avec les mêmes étapes, toutes décochées. \\ \hline
Supprimer & Supprime la tâche. Un message propose \btn{Annuler} pendant quelques secondes. \\ \hline
Fermer & La croix, ou la touche \kbd{Échap}. \\ \hline
\end{lexique}

\subsection{Reporter une tâche bloquée}
Le client n’a pas envoyé ses pièces à temps ? Tu reportes la tâche et tu prends autre chose à la place.
\begin{etapes}
  \item Ouvre la tâche, puis clique sur \btn{Reporter} en haut du panneau.
  \item Choisis la raison : pièces manquantes du client, attente du manager, ou autre.
  \item Choisis la date de reprise : \btn{Demain}, \btn{Dans 3 jours ouvrés}, \btn{Lundi prochain}, \btn{Dans 2 semaines}, ou une date au choix.
  \item Si besoin, change l’échéance et précise ce qui manque (« relevés bancaires d’octobre »).
  \item Laisse cochées les deux options : la tâche passe \textbf{En attente client}, et les créneaux prévus pour ce dossier d’ici la reprise sont \textbf{libérés} dans la répartition.
  \item Clique sur \btn{Reporter au …}.
\end{etapes}
\capture[0.8\linewidth]{19-reporter.jpg}{La fenêtre « Reporter la tâche ».}
\begin{attention}
Si la reprise tombe après l’échéance, Nini l’affiche en orange : préviens le manager du dossier, ou repousse l’échéance si c’est possible.
Les échéances légales (TVA…) ne se repoussent pas.
\end{attention}

\textbf{Ensuite, Nini te propose quoi faire à la place} :
\begin{itemize}
  \item \textbf{Disponibles maintenant} : les tâches des autres dossiers que tu peux traiter tout de suite, les plus urgentes d’abord ;
  \item \textbf{À avancer} : les tâches des périodes suivantes que tu peux commencer en avance.
\end{itemize}
Les autres tâches du dossier bloqué ne sont pas proposées : elles attendent sans doute les mêmes pièces.
Clique sur \btn{Prendre} : si du temps a été libéré aujourd’hui, il est donné à ce dossier dans ta répartition, et la tâche s’ouvre.
\capture[0.8\linewidth]{20-a-la-place.jpg}{« Que faire à la place ? » : un clic sur Prendre.}
\begin{astuce}
Le jour de la reprise, la tâche réapparaît en haut de l’accueil, dans \btn{À reprendre}.
Tu peux aussi la reprendre plus tôt : ouvre-la et clique sur \btn{Reprendre maintenant} dans le bandeau orange.
Les jours libérés restent libres dans la répartition : tu peux y ajouter un créneau avec \btn{+ Ajouter}.
\end{astuce}

\subsection{Créer une tâche}
Trois façons :
\begin{itemize}
  \item le bouton bleu \btn{Nouvelle tâche}, ou la touche \kbd{N} : écris le titre, la tâche est créée ;
  \item le champ \btn{Ajouter une tâche…} en haut des listes : écris le titre puis \kbd{Entrée}. À droite du champ, la liste \btn{Lier à : …} rattache directement la tâche à un dossier ;
  \item dans le calendrier : clique sur un jour.
\end{itemize}

% =====================================================================
\section{Toutes les tâches}
% =====================================================================
\capture{04-taches.jpg}{La liste de toutes les tâches, avec les filtres en haut.}
Les filtres se combinent : dossier, statut, échéance (en retard, à traiter en ce moment, cette semaine…), catégorie, priorité.
Un filtre actif s’affiche en bleu. \btn{Effacer les filtres} remet tout à zéro. Le champ \btn{Filtrer…} cherche dans les titres, les notes et les noms de dossiers.

\subsection{Regrouper les tâches}
La liste \btn{Regrouper…} range les tâches :
\begin{lexique}
Par dossier & Un groupe par client (choix par défaut). Les groupes sont repliés : clique sur un nom pour voir ses tâches. \\ \hline
Par manager & Un groupe par manager, avec sa couleur. \\ \hline
Par période du mois & Début, milieu, TVA, fin de mois, puis les tâches hors routine. \\ \hline
Par statut, catégorie, priorité & Pour voir d’un coup tout ce qui est en attente, toutes les TVA, toutes les urgences… \\ \hline
Sans regroupement & Une seule liste, triée selon le tri choisi. \\ \hline
\end{lexique}
\btn{Tout déplier} et \btn{Tout replier} ouvrent ou ferment tous les groupes. Le nombre de tâches en retard s’affiche en rouge sur chaque groupe.

\subsection{Lier une tâche à un dossier}
\textbf{Dès la création} : dans le champ \btn{Ajouter une tâche…}, choisis le dossier dans la liste à droite (\btn{Lier à : …}), puis écris le titre.
\capture[0.8\linewidth]{22-ajout-lie.jpg}{Créer une tâche déjà liée à un dossier.}
\textbf{Plus tard} : une tâche créée sans dossier apparaît dans le groupe \btn{Sans dossier — à lier}.
Ouvre-la, puis choisis le dossier dans \btn{Dossier}. Elle rejoint aussitôt le groupe de ce dossier.
\capture[0.8\linewidth]{21-a-lier.jpg}{Le groupe « Sans dossier — à lier ». Les tâches générales du cabinet ont leur propre groupe.}
\begin{astuce}
Les \textbf{tâches générales} de la routine (tableau de variation du CA, point avec chaque manager) ne concernent pas un seul dossier : elles sont rangées à part, dans \btn{Tâches générales (cabinet)}.
\end{astuce}

\subsection{La vue Tableau}
\capture{05-tableau.jpg}{La vue Tableau : une colonne par statut.}
Clique sur \btn{Tableau} en haut à droite. Pour changer le statut d’une tâche, \textbf{fais glisser sa carte} dans une autre colonne.
La colonne « Terminé » montre les 14 derniers jours.

% =====================================================================
\section{Organiser le mois : la routine et la répartition}
% =====================================================================
\subsection{Le mois type}
Chaque dossier passe par quatre périodes. Nini crée une tâche par dossier et par période, avec ses étapes :

\begin{tikzpicture}[x=0.995\linewidth/31]
  \def\h{3.05}
  \fill[accentsoft, rounded corners=6pt] (0,0) rectangle (10,\h);
  \fill[white] (10,0) rectangle (17,\h);
  \fill[ambersoft] (17,0) rectangle (24,\h);
  \fill[white, rounded corners=6pt] (24,0) rectangle (31,\h);
  \draw[line, rounded corners=6pt, line width=0.6pt] (0,0) rectangle (31,\h);
  \foreach \x in {10,17,24} \draw[line, line width=0.6pt] (\x,0) -- (\x,\h);
  \newcommand{\periode}[5]{%
    \node[anchor=north west, text width=#2*\linewidth/31-14pt, align=left, inner sep=0pt] at (#1+0.35,\h-0.18) {%
      {\ttfamily\small\color{ink2}#3}\\[2pt]\textbf{#4}\\[2pt]{\footnotesize\color{ink2}#5}};}
  \periode{0}{10}{1 → 10}{Début de mois}{Banque, recettes, lettrage, 471, balance. À finir avant le 17.}
  \periode{10}{7}{11 → 17}{Milieu}{Achats, paie, OD, états manquants.}
  \periode{17}{7}{18 → 24}{TVA}{Déclarations du 19, du 21 et du 24. Révision.}
  \periode{24}{7}{25 → fin}{Fin de mois}{Tableaux de bord, finalisation, point manager.}
  \fill[pink, rounded corners=1.5pt] (4,-0.5) rectangle (10,-0.36);
  \node[anchor=west, font=\footnotesize, text=ink2] at (10.3,-0.43) {Du 5 au 10 : semaine réservée au dossier qui a une période réservée.};
\end{tikzpicture}

En fin de mois, Nini crée aussi le \textbf{tableau de variation du CA} et un \textbf{point dossiers} pour chaque manager, avec la liste de ses dossiers.

\subsection{Le 1\textsuperscript{er} jour de chaque mois}
\begin{etapes}
  \item Ouvre \btn{Répartition}.
  \item Clique sur \btn{Créer les tâches du mois}. Si tu cliques deux fois, rien n’est créé en double.
  \item Clique sur \btn{Proposer une répartition}.
  \item Regarde le planning et ajuste ce qui ne te convient pas.
\end{etapes}

\capture{06-repartition.jpg}{La répartition du mois : un bloc par dossier et par jour, l’équilibre à droite.}

\subsection{Lire le planning}
\begin{itemize}
  \item \textbf{Une case = un jour ouvré}. Les jours fériés sont grisés.
  \item En haut de chaque case : la \textbf{période} et le temps total du jour (« 1 j », ou « 0,25 j libre »). En rouge si la journée dépasse 1 jour.
  \item \textbf{Chaque bloc} : le code du dossier, son nom et la durée. La couleur du bord indique le manager.
  \item En haut de la page : le nombre de jours ouvrés du mois, le temps alloué et la marge.
\end{itemize}

\subsection{Modifier le planning}
\begin{itemize}
  \item \textbf{Déplacer} : fais glisser un bloc vers un autre jour.
  \item \textbf{Changer} : clique sur un bloc pour choisir un autre dossier, un autre jour, une autre durée ou une autre période. Le bouton \btn{Supprimer} est en bas à gauche.
  \item \textbf{Ajouter} : passe la souris sur un jour qui n’est pas plein, puis clique sur \btn{+ Ajouter}.
  \item \textbf{Tout recalculer} : \btn{Nouvelle proposition}. Attention, cela remplace tes modifications du mois.
  \item \textbf{Autre mois} : les flèches en haut à droite.
\end{itemize}
\capture[0.7\linewidth]{07-creneau.jpg}{La fenêtre d’un créneau : dossier, jour, durée, période.}

\subsection{L’équilibre par dossier}
À droite du planning, chaque dossier montre son \textbf{temps alloué} par mois. Tu peux le modifier directement dans la case.
À côté, l’\textbf{écart} : \textcolor{ok}{\symb ✓} si le planning correspond, \textcolor{amber}{+} s’il y a trop de temps prévu, \textcolor{no}{−} s’il en manque.

\subsection{Comment Nini propose la répartition}
\begin{itemize}
  \item Le dossier qui a une \textbf{période réservée} (par exemple du 5 au 10) est placé en premier sur ces jours. Si un jour est férié, il continue le jour ouvré suivant.
  \item Le temps de chaque autre dossier est \textbf{réparti sur les quatre périodes}, pour que chaque dossier soit traité à chaque étape du mois.
  \item Les petits dossiers passent d’abord en début de mois (la banque), puis en TVA.
  \item En période TVA, les dossiers dont la date limite est au \textbf{19} passent avant ceux au 21, puis au 24. Un dossier sans TVA n’y est jamais placé.
  \item Les dossiers d’un même manager sont regroupés.
\end{itemize}

\subsection{Modifier la routine}
Les périodes, leurs dates et leurs étapes se modifient dans \btn{Réglages → Routine mensuelle} (voir la partie \ref{sec:reglages}).

% =====================================================================
\section{Les dossiers}
% =====================================================================
\capture{08-dossiers.jpg}{La liste des dossiers.}
Pour chaque dossier : le manager, le temps alloué, l’avancement des échéances de l’année, les tâches ouvertes et en retard, la prochaine échéance et le temps passé.
Coche \btn{Afficher les archivés} pour voir aussi les dossiers archivés.

\subsection{Regrouper les dossiers}
La liste \btn{Regrouper par…} range les dossiers par \textbf{manager} (choix par défaut), par \textbf{régime de TVA}, par \textbf{forme juridique} ou par \textbf{imposition}, ou sans regroupement.
Clique sur l’en-tête d’un groupe pour le replier ou le déplier : il affiche le nombre de dossiers, le temps alloué par mois et les retards.
\btn{Tout déplier} et \btn{Tout replier} agissent sur tous les groupes. Nini se souvient de ton choix.

\subsection{La fiche d’un dossier}
\capture{09-fiche.jpg}{La fiche d’un dossier : chiffres, tâches, fiche et notes.}
\begin{itemize}
  \item En haut : tâches ouvertes, retards, attente client, échéances traitées, temps passé par rapport au budget.
  \item Au milieu : les tâches \btn{À faire} et \btn{Terminées}, avec un champ pour en ajouter.
  \item À droite : la fiche (le crayon sert à la modifier), puis les \textbf{notes} du dossier, enregistrées automatiquement.
  \item Boutons du haut : \btn{Poser les échéances}, \btn{Rapport} (le rapport de ce dossier), \btn{Tâche}.
  \item En bas à droite : \btn{Archiver} (le dossier disparaît des listes, sans être effacé) ou \btn{Supprimer} (définitif, avec ses tâches).
\end{itemize}

\subsection{Créer ou modifier un dossier}
\capture[0.75\linewidth]{10-formulaire.jpg}{Le formulaire d’un dossier.}
\begin{lexique}
Nom et code & Le code (8 lettres maximum) s’affiche partout. S’il est vide, Nini en propose un. \\ \hline
Manager & Choisis-le dans la liste, ou tape un nouveau prénom. \\ \hline
Temps alloué & En jours par mois, par exemple \texttt{0.5} ou \texttt{1.25}. Il sert à la répartition. \\ \hline
Période réservée & À cocher si le dossier a des jours réservés (du 5 au 10, par exemple). \\ \hline
Imposition & IS ou IR. \\ \hline
Clôture & La date de fin d’exercice, par exemple \texttt{31/12}. \\ \hline
Régime de TVA & Mensuel, trimestriel, CA12 ou sans TVA. Sans TVA, le dossier n’a pas d’étapes TVA. \\ \hline
Date limite TVA & Le 19, le 21 ou le 24. Elle sert d’échéance à la tâche TVA. \\ \hline
Paie & À cocher si le cabinet gère la paie. \\ \hline
Contact, budget & Interlocuteur, e-mail, téléphone, honoraires, budget temps annuel. \\ \hline
\end{lexique}
Quand tu crées un dossier, Nini te propose tout de suite de poser ses échéances de l’année.

% =====================================================================
\section{Poser les échéances de l’année}
% =====================================================================
\capture[0.8\linewidth]{11-echeances.jpg}{Les échéances proposées pour l’année.}
Ce bouton (dans un dossier, ou dans la liste des dossiers pour tous à la fois) propose les grandes échéances de l’année :
acomptes et solde d’IS, liasse, AG, dépôt au greffe, CFE, paie, TVA.
\begin{etapes}
  \item Choisis l’année.
  \item Coche ou décoche chaque échéance. Celles déjà passées sont décochées, celles déjà créées sont grisées.
  \item Clique sur \btn{Créer}.
\end{etapes}
\begin{attention}
La routine du mois crée déjà la TVA mensuelle. Ici, \textbf{décoche les déclarations de TVA} pour ne pas les avoir en double.
Les dates proposées sont des repères : la règle du cabinet prime.
\end{attention}

% =====================================================================
\section{Le calendrier}
% =====================================================================
\capture{12-calendrier.jpg}{Le calendrier des échéances.}
Chaque tâche apparaît le jour de son échéance, avec le code du dossier. En rouge : en retard. Barré : terminé.
Clique sur une tâche pour l’ouvrir, ou sur un jour vide pour créer une tâche à cette date.
La liste en haut à droite filtre par dossier.

% =====================================================================
\section{Les rapports}
% =====================================================================
\capture{13-rapports.jpg}{Le rapport d’activité du mois.}
\begin{etapes}
  \item Choisis la période : cette semaine, ce mois, le mois dernier, le trimestre, l’année, ou des dates libres.
  \item Choisis un dossier, ou garde « Tous les dossiers ».
  \item Options : \btn{Inclure les échéances futures} et \btn{Temps passés}.
\end{etapes}
Le rapport montre une synthèse par dossier, puis le détail : \textbf{travaux réalisés} et \textbf{reste à faire}, avec les retards.
\begin{itemize}
  \item \btn{Imprimer ou PDF} : imprime le rapport, ou choisis « Enregistrer au format PDF » dans la fenêtre d’impression.
  \item \btn{Excel (CSV)} : télécharge le détail pour Excel.
\end{itemize}

% =====================================================================
\section{Les réglages}\label{sec:reglages}
% =====================================================================
\subsection{Profil}
Ton prénom, affiché sur l’accueil et sur les rapports.

\subsection{La routine mensuelle}
\capture{15-routine.jpg}{La routine : une carte par période.}
Pour chaque période :
\begin{lexique}
Nom & Le nom de la période (« Début de mois »…). \\ \hline
Planifiée du … au … & Les jours du mois où ce travail est prévu dans la répartition. C’est aussi la date de début des tâches. \\ \hline
Échéance & Un jour fixe (le 17…), la date TVA du dossier, ou le dernier jour ouvré. \\ \hline
Tâches & Une par dossier, une par manager, ou une seule. \\ \hline
Étapes & Une étape par ligne. \\ \hline
\end{lexique}
\btn{Ajouter une période}, la corbeille pour en supprimer une, \btn{Routine d’origine} pour revenir au départ, puis \btn{Enregistrer la routine}.
Les changements s’appliquent aux tâches créées ensuite, pas à celles qui existent déjà.

\subsection{Stockage, sauvegardes et exports}
\capture{16-sauvegardes.jpg}{Stockage, sauvegardes, exports et données.}
\begin{lexique}
Exporter une sauvegarde & Télécharge un fichier \texttt{.json} avec toutes tes données. À faire chaque vendredi, et à ranger sur le lecteur du cabinet. \\ \hline
Restaurer… & Remplace toutes les données par celles d’un fichier de sauvegarde. \\ \hline
Exports Excel & Toutes les tâches, ou tous les dossiers, en fichier pour Excel. \\ \hline
Données d’exemple & Remplace tes données par des dossiers d’exemple, pour t’entraîner. \\ \hline
Tout effacer & Efface tout. Nini demande une confirmation. \\ \hline
\end{lexique}

% =====================================================================
\section{La sécurité}
% =====================================================================
\begin{securite}
\textbf{Nini ne peut rien envoyer sur Internet.} Le navigateur lui interdit toute connexion. Même si quelqu’un essayait, la demande serait bloquée avant de partir.
\end{securite}

\capture{14-securite.jpg}{Réglages → Sécurité : l’état du verrou, le test et le journal.}

\subsection{Vérifier que rien ne sort}
\begin{etapes}
  \item Va dans \btn{Réglages → Sécurité}.
  \item Le bandeau doit être \textbf{vert} : « Verrou réseau actif ».
  \item Clique sur \btn{Tester le verrou} : Nini essaie d’envoyer des données de 4 façons. Les 4 doivent être \textbf{bloquées}.
  \item Le compteur \btn{requête vers l’extérieur} doit rester à \textbf{0}. Sinon, le bandeau devient rouge.
\end{etapes}

\subsection{Le journal de sécurité}
Nini note tout ce qui concerne tes données, avec la date et l’heure : chaque ouverture, chaque tentative bloquée, chaque export et chaque restauration.
Le journal ne peut pas être effacé depuis l’app. \btn{Exporter le journal} donne un fichier Excel à garder ou à montrer.

\subsection{Vérifier sans faire confiance à Nini}
\begin{itemize}
  \item \textbf{Sur le PC} : touche \kbd{F12}, onglet \btn{Réseau}. Aucune ligne ne doit apparaître quand tu utilises Nini.
  \item \textbf{Sur l’iPad} : passe en mode avion. Nini fonctionne exactement pareil.
\end{itemize}

\subsection{Les règles de sécurité}
\begin{lexique}[0.42]
\textcolor{ok}{\symb ✓}\enspace Sauvegarde chaque vendredi & Sur le lecteur du cabinet. Nini te le rappelle après 7 jours. \\ \hline
\textcolor{ok}{\symb ✓}\enspace Verrouille ton ordinateur & \kbd{Windows}+\kbd{L} quand tu t’absentes : Nini n’a pas de mot de passe, il est protégé par ta session. \\ \hline
\textcolor{ok}{\symb ✓}\enspace Ouvre Nini de la même façon & Même fichier, même endroit, même navigateur. \\ \hline
\textcolor{no}{\symb ✕}\enspace Ne vide pas le navigateur & Effacer les données de navigation efface Nini. \\ \hline
\textcolor{no}{\symb ✕}\enspace Pas de navigation privée & Tout serait perdu à la fermeture. \\ \hline
\textcolor{no}{\symb ✕}\enspace Pas de mail perso ni de cloud perso & Une sauvegarde contient les noms des clients : c’est un document confidentiel du cabinet. \\ \hline
\end{lexique}

% =====================================================================
\section{L’iPad et le changement d’appareil}\label{sec:ipad}
% =====================================================================
\subsection{Installer Nini sur l’iPad (une seule fois)}
\begin{etapes}
  \item Avec Internet, ouvre l’adresse de Nini dans \btn{Safari}.
  \item Touche \btn{Partager → Sur l’écran d’accueil → Ajouter}. L’icône bleue « N » apparaît.
  \item Ouvre Nini \btn{depuis son icône}, pas depuis Safari.
  \item \btn{Réglages → Restaurer} : choisis ta sauvegarde, reçue par AirDrop (elle est dans l’app Fichiers).
  \item C’est fini : coupe le Wi-Fi ou passe en mode avion, Nini fonctionne.
\end{etapes}
\capture[0.55\linewidth]{18-ipad.jpg}{Nini sur iPad, en mode portrait. Le menu s’ouvre avec le bouton en haut à gauche.}
\begin{attention}
Ne supprime pas l’icône de Nini : cela effacerait les données sur l’iPad.
Sur l’iPad, \btn{Exporter une sauvegarde} ouvre le menu de partage : choisis \btn{AirDrop} ou \btn{Enregistrer dans Fichiers}.
\end{attention}

\subsection{Passer d’un appareil à l’autre}
\begin{etapes}
  \item Sur l’appareil où tu as travaillé : \btn{Exporter une sauvegarde}.
  \item Transporte le fichier : AirDrop, câble, clé USB ou lecteur du cabinet.
  \item Sur l’autre appareil : \btn{Réglages → Restaurer}.
\end{etapes}
\begin{attention}
Ne travaille jamais sur deux appareils en même temps : la restauration remplace tout ce qu’il y avait avant.
\end{attention}

% =====================================================================
\section{Raccourcis et palette de recherche}
% =====================================================================
\capture[0.8\linewidth]{17-palette.jpg}{La palette : tape quelques lettres, puis Entrée.}
\begin{lexique}[0.32]
\kbd{Ctrl} + \kbd{K} & Ouvre la palette : chercher une tâche, un dossier, aller à une page, créer, lancer ou arrêter un chrono. Sur Mac : \kbd{⌘}+\kbd{K}. \\ \hline
\kbd{N} & Nouvelle tâche (quand tu n’es pas en train d’écrire). \\ \hline
\kbd{Échap} & Ferme le panneau ou la fenêtre ouverte. \\ \hline
\kbd{Entrée} & Valide un champ : étape, temps, titre, ajout rapide. \\ \hline
\kbd{Ctrl} + \kbd{Entrée} & Publie un commentaire. \\ \hline
\end{lexique}

% =====================================================================
\section{Questions fréquentes}
% =====================================================================
\begin{lexique}[0.36]
Je ne vois plus mes dossiers & \btn{Réglages → Restaurer}, puis choisis ta dernière sauvegarde. Vérifie aussi que tu as ouvert le bon fichier, dans le bon navigateur. \\ \hline
Le client n’a pas envoyé ses pièces & Ouvre la tâche → \btn{Reporter}, choisis une date de reprise, puis prends une autre tâche à la place. \\ \hline
J’ai coché une tâche par erreur & \btn{Annuler} dans le message en bas, ou clique à nouveau sur le rond. \\ \hline
J’ai supprimé une tâche par erreur & \btn{Annuler} dans le message en bas, tout de suite. Sinon, restaure ta dernière sauvegarde. \\ \hline
Les tâches du mois n’apparaissent pas & \btn{Répartition → Créer les tâches du mois}. \\ \hline
Le planning ne me convient pas & Déplace les blocs, ou \btn{Nouvelle proposition}. \\ \hline
Un dossier a trop ou pas assez de temps & Change son temps alloué dans l’équilibre (à droite du planning), puis \btn{Nouvelle proposition}. \\ \hline
La date de TVA d’un dossier change & Fiche du dossier → crayon → \btn{Date limite TVA}. \\ \hline
J’ai créé une tâche sans dossier & Ouvre-la, puis choisis le dossier dans \btn{Dossier}. Les tâches à lier sont regroupées dans \btn{Sans dossier — à lier}. \\ \hline
Je veux voir mes dossiers autrement & \btn{Dossiers} → \btn{Regrouper par…} : manager, TVA, forme ou imposition. \\ \hline
Un nouveau client arrive & \btn{Dossiers → Nouveau dossier}, puis pose ses échéances. \\ \hline
Un client part & Fiche du dossier → \btn{Archiver}. Son historique est gardé. \\ \hline
Je veux changer une étape pour tous les mois & \btn{Réglages → Routine mensuelle}. \\ \hline
Le rapport ne montre pas tout & Coche \btn{Inclure les échéances futures}, ou choisis une période plus longue. \\ \hline
Le cadre jaune de sauvegarde reste affiché & Clique sur \btn{Sauvegarder} : il disparaît pour 7 jours. \\ \hline
\end{lexique}

% =====================================================================
\section*{Les règles d’or}
% =====================================================================
\addcontentsline{toc}{section}{Les règles d’or}
\begin{securite}
\begin{etapes}
  \item Le 1\textsuperscript{er} du mois : \btn{Créer les tâches du mois}, puis \btn{Proposer une répartition}.
  \item Chaque matin : \btn{Aujourd’hui}, puis coche les étapes au fil de la journée.
  \item Chaque vendredi : \btn{Exporter une sauvegarde} sur le lecteur du cabinet.
  \item Toujours : verrouiller l’ordinateur, pas de navigation privée, pas de sauvegarde par mail perso.
  \item En cas de doute sur une échéance : la règle du cabinet prime.
\end{etapes}
\end{securite}

\end{document}
